\documentclass[11pt,a4paper]{article}

%==============================================================================
% PACKAGES
%==============================================================================
\usepackage[utf8]{inputenc}
\usepackage[T1]{fontenc}
\usepackage[french]{babel}
\usepackage{geometry}
\usepackage{graphicx}
\usepackage{xcolor}
\usepackage{hyperref}
\usepackage{fancyhdr}
\usepackage{titlesec}
\usepackage{enumitem}
\usepackage{tcolorbox}
\usepackage{tabularx}
\usepackage{booktabs}
\usepackage{fontawesome5}
\usepackage{tikz}
\usepackage{multicol}
\usepackage{float}
\usepackage{etoolbox}
\usepackage{setspace}
\usepackage{soul}
\usepackage{mdframed}

\usetikzlibrary{shadows.blur, positioning, arrows.meta, shapes.geometric, calc, decorations.pathreplacing}

%==============================================================================
% GÉOMÉTRIE ET COULEURS
%==============================================================================
\geometry{
    margin=2cm,
    top=2.5cm,
    bottom=2cm,
    headheight=15pt
}

% Palette de couleurs moderne
\definecolor{primary}{RGB}{41, 65, 114}
\definecolor{secondary}{RGB}{52, 152, 219}
\definecolor{accent}{RGB}{231, 76, 60}
\definecolor{success}{RGB}{39, 174, 96}
\definecolor{warning}{RGB}{243, 156, 18}
\definecolor{dark}{RGB}{44, 62, 80}
\definecolor{light}{RGB}{236, 240, 241}
\definecolor{purple}{RGB}{155, 89, 182}

%==============================================================================
% CONFIGURATION DES LIENS
%==============================================================================
\hypersetup{
    colorlinks=true,
    linkcolor=primary,
    urlcolor=secondary,
    pdfauthor={Projet Binôme},
    pdftitle={Préparation Soutenance - Simulateur de Parking},
}

%==============================================================================
% STYLE DES TITRES
%==============================================================================
\titleformat{\part}[display]
    {\centering\Huge\bfseries\color{primary}}
    {\textcolor{secondary}{Partie \thepart}}
    {0.5em}
    {}
    [\vspace{0.5em}\textcolor{primary}{\rule{0.5\linewidth}{2pt}}]

\titleformat{\section}
    {\Large\bfseries\color{primary}}
    {\tikz[baseline=(char.base)]{
        \node[fill=primary, text=white, rounded corners=3pt, inner sep=5pt] (char) {\thesection};
    }}
    {0.8em}
    {}

\titleformat{\subsection}
    {\large\bfseries\color{dark}}
    {\textcolor{secondary}{\faIcon{chevron-right}}}
    {0.5em}
    {}

%==============================================================================
% EN-TÊTE ET PIED DE PAGE
%==============================================================================
\pagestyle{fancy}
\fancyhf{}
\renewcommand{\headrulewidth}{0.5pt}
\renewcommand{\footrulewidth}{0.3pt}
\fancyhead[L]{\textcolor{primary}{\faIcon{microphone-alt} \textbf{Préparation Soutenance}}}
\fancyhead[R]{\textcolor{dark}{Simulateur de Parking}}
\fancyfoot[L]{\textcolor{dark}{\small Durée : 10 minutes}}
\fancyfoot[C]{\textcolor{primary}{\thepage}}
\fancyfoot[R]{\textcolor{dark}{\small 2024/2025}}

%==============================================================================
% BOÎTES PERSONNALISÉES
%==============================================================================
\tcbuselibrary{skins, breakable}

% Boîte pour le timing
\newtcolorbox{timingbox}[1][]{
    enhanced,
    colback=secondary!5,
    colframe=secondary,
    arc=5pt,
    boxrule=2pt,
    left=10pt,
    right=10pt,
    top=5pt,
    bottom=5pt,
    fonttitle=\bfseries\large\color{white},
    attach boxed title to top center={yshift=-3mm},
    boxed title style={colback=secondary, arc=3pt, boxrule=0pt},
    title={\faIcon{clock} #1}
}

% Boîte pour les questions
\newtcolorbox{questionbox}[1][]{
    enhanced,
    breakable,
    colback=primary!5,
    colframe=primary,
    arc=4pt,
    boxrule=1.5pt,
    left=12pt,
    right=12pt,
    top=10pt,
    bottom=10pt,
    fonttitle=\bfseries\color{white},
    attach boxed title to top left={yshift=-3mm, xshift=8mm},
    boxed title style={colback=primary, arc=3pt},
    title={\faIcon{question-circle} #1}
}

% Boîte pour les réponses
\newtcolorbox{reponsebox}{
    enhanced,
    colback=success!5,
    colframe=success!70!black,
    arc=4pt,
    boxrule=1pt,
    left=12pt,
    right=12pt,
    top=8pt,
    bottom=8pt,
    before upper={\textcolor{success!70!black}{\faIcon{comment-dots} \textbf{Réponse :}} }
}

% Boîte conseil
\newtcolorbox{conseilbox}[1][]{
    enhanced,
    colback=warning!8,
    colframe=warning,
    arc=4pt,
    boxrule=1pt,
    left=10pt,
    right=10pt,
    fonttitle=\bfseries\color{white},
    attach boxed title to top left={yshift=-3mm, xshift=5mm},
    boxed title style={colback=warning, arc=2pt},
    title={\faIcon{lightbulb} #1}
}

% Boîte pour les idées clés
\newtcolorbox{ideesbox}{
    enhanced,
    colback=light,
    colframe=dark!50,
    arc=3pt,
    boxrule=0.5pt,
    left=8pt,
    right=8pt,
    top=5pt,
    bottom=5pt,
    before upper={\textcolor{dark}{\faIcon{key} \textbf{Idées clés :}}\par\smallskip}
}

% Boîte phrase clé
\newtcolorbox{phrasebox}{
    enhanced,
    colback=purple!8,
    colframe=purple,
    arc=3pt,
    boxrule=1pt,
    left=10pt,
    right=10pt,
    before upper={\textcolor{purple}{\faIcon{quote-left}} }
}

%==============================================================================
% COMMANDES PERSONNALISÉES
%==============================================================================
\newcommand{\keyword}[1]{\textcolor{primary}{\textbf{#1}}}
\newcommand{\code}[1]{\texttt{\textcolor{accent}{#1}}}
\newcommand{\timing}[1]{\textcolor{secondary}{\faIcon{stopwatch} \textbf{#1}}}
\newcommand{\important}[1]{\textcolor{accent}{\textbf{#1}}}

% Numérotation des questions
\newcounter{questioncounter}
\setcounter{questioncounter}{0}
\newcommand{\question}[1]{%
    \stepcounter{questioncounter}%
    \begin{questionbox}[Q\thequestioncounter\ : #1]%
    \end{questionbox}%
}

%==============================================================================
% DOCUMENT
%==============================================================================
\begin{document}

%------------------------------------------------------------------------------
% PAGE DE TITRE
%------------------------------------------------------------------------------
\begin{titlepage}
    \begin{tikzpicture}[remember picture, overlay]
        % Fond coloré en haut
        \fill[primary] (current page.north west) rectangle ([yshift=-7cm]current page.north east);

        % Cercles décoratifs
        \foreach \i in {1,...,15} {
            \fill[secondary, opacity=0.1] ($(current page.north west) + (\i*1.2, -\i*0.4)$) circle (1.5cm);
        }

        % Barre
        \fill[secondary] ([yshift=-7cm]current page.north west) rectangle ([yshift=-7.4cm]current page.north east);

        % Icône microphone
        \node[anchor=center] at ([yshift=-3.5cm]current page.north) {
            \textcolor{white}{\fontsize{70}{70}\selectfont\faIcon{microphone-alt}}
        };
    \end{tikzpicture}

    \vspace*{5.5cm}

    \begin{center}
        {\Huge\bfseries\textcolor{primary}{Préparation à la Soutenance}}\\[0.8cm]
        {\LARGE\textcolor{dark}{Simulateur de Parking en C}}\\[1.5cm]

        \textcolor{secondary}{\rule{10cm}{1.5pt}}\\[1.5cm]

        \begin{tikzpicture}
            \node[fill=accent!20, rounded corners=8pt, inner sep=15pt] {
                \begin{minipage}{8cm}
                    \centering
                    {\Large\textcolor{accent}{\faIcon{clock} Durée totale : 10 minutes}}\\[0.5cm]
                    \textcolor{dark}{
                        \faIcon{users} Présentation : 5 min\\
                        \faIcon{comments} Questions : 5 min
                    }
                \end{minipage}
            };
        \end{tikzpicture}

        \vspace{2cm}

        \begin{tcolorbox}[
            width=12cm,
            colback=light,
            colframe=primary,
            arc=5pt,
            boxrule=1pt,
            halign=center
        ]
            {\large\textbf{Ce document n'est pas à lire pendant la soutenance}}\\[0.3cm]
            \textcolor{dark}{Il sert à \textbf{préparer} et \textbf{mémoriser} les points clés}
        \end{tcolorbox}

        \vfill

        \textcolor{dark}{\faIcon{graduation-cap} Formation Ingénieur - 2024/2025}
    \end{center}
\end{titlepage}

%------------------------------------------------------------------------------
% TABLE DES MATIÈRES
%------------------------------------------------------------------------------
\newpage
\thispagestyle{empty}

\begin{tikzpicture}[remember picture, overlay]
    \fill[secondary!15] (current page.north west) rectangle ([xshift=0.6cm]current page.south west);
\end{tikzpicture}

\tableofcontents
\newpage

%##############################################################################
\part{Plan de Présentation}
\addcontentsline{toc}{part}{Partie I : Plan de Présentation (5 minutes)}
%##############################################################################

\begin{conseilbox}[Conseils pour l'oral]
\begin{multicols}{2}
\begin{itemize}[noitemsep]
    \item[\faIcon{eye}] Regarder le jury, pas l'écran
    \item[\faIcon{comment}] Parler clairement, pas trop vite
    \item[\faIcon{file-alt}] Ne pas lire ses notes
    \item[\faIcon{laptop}] Préparer une démo fonctionnelle
    \item[\faIcon{smile}] Rester détendu et naturel
    \item[\faIcon{users}] S'adresser à tous les jurys
\end{itemize}
\end{multicols}
\end{conseilbox}

\vspace{0.5cm}

%==============================================================================
\section{Introduction}
%==============================================================================

\begin{timingbox}[30 secondes]
\end{timingbox}

\begin{ideesbox}
\begin{itemize}[noitemsep]
    \item Présentation rapide du binôme
    \item Objectif : simulateur de parking en C avec ncurses
    \item Le joueur contrôle les barrières, les voitures sont autonomes
    \item Collision = Game Over
\end{itemize}
\end{ideesbox}

\begin{phrasebox}
"Bonjour, nous allons vous présenter notre simulateur de parking développé en C. L'objectif était de créer un jeu où des véhicules autonomes circulent dans un parking, et où le joueur doit gérer l'ouverture des barrières pour éviter les collisions."
\end{phrasebox}

\textbf{Transition :} \textit{"Commençons par voir comment fonctionne le jeu en pratique..."}

%==============================================================================
\section{Principe du jeu}
%==============================================================================

\begin{timingbox}[1 minute]
\end{timingbox}

\begin{ideesbox}
\begin{itemize}[noitemsep]
    \item Affichage console avec ncurses (couleurs, UTF-8)
    \item Véhicules dans une file d'attente à l'entrée
    \item Touches E et S pour les barrières
    \item Les voitures suivent des flèches et trouvent une place
    \item Après un temps, elles repartent vers la sortie
    \item Deux modes : Normal et Hard
\end{itemize}
\end{ideesbox}

\begin{conseilbox}[Démonstration suggérée]
\begin{enumerate}[noitemsep]
    \item Lancer le jeu, montrer le menu de difficulté
    \item Ouvrir la barrière d'entrée, laisser un véhicule entrer
    \item Montrer le parking automatique d'un véhicule
    \item Attendre qu'il reparte et ouvrir la sortie
    \item Montrer le HUD (score, places, file d'attente)
\end{enumerate}
\end{conseilbox}

\textbf{Transition :} \textit{"Voyons maintenant comment tout ça fonctionne en interne..."}

%==============================================================================
\section{Fonctionnement interne}
%==============================================================================

\begin{timingbox}[2 minutes]
\end{timingbox}

\subsection{Structure du code \timing{30 sec}}

\begin{ideesbox}
\begin{itemize}[noitemsep]
    \item Projet découpé en modules : \code{jeu.c}, \code{mouvement.c}, \code{affichage.c}, \code{plan.c}
    \item Chaque module a un rôle précis (séparation des responsabilités)
    \item Véhicules stockés dans une liste chaînée
\end{itemize}
\end{ideesbox}

\subsection{Gestion des véhicules \timing{45 sec}}

\begin{ideesbox}
\begin{itemize}[noitemsep]
    \item Un véhicule = structure avec position, direction, état, sprite
    \item Cycle de vie : création $\rightarrow$ file $\rightarrow$ entrée $\rightarrow$ navigation $\rightarrow$ parking $\rightarrow$ sortie
    \item État '1' = actif, '0' = garé
\end{itemize}
\end{ideesbox}

\subsection{Déplacement \timing{45 sec}}

\begin{ideesbox}
\begin{itemize}[noitemsep]
    \item Les véhicules suivent les flèches du plan
    \item Système de \keyword{scoring} : on calcule un score pour chaque direction
    \item La direction avec le meilleur score est choisie
    \item Vérification qu'il n'y a pas de véhicule devant
\end{itemize}
\end{ideesbox}

\textbf{Transition :} \textit{"Quelques points techniques importants..."}

%==============================================================================
\section{Points techniques}
%==============================================================================

\begin{timingbox}[1 minute]
\end{timingbox}

\subsection{Stabilité des déplacements \timing{30 sec}}

\begin{ideesbox}
\begin{itemize}[noitemsep]
    \item \textbf{Problème} : les voitures oscillaient entre deux directions
    \item \textbf{Solution} : \keyword{lane-lock} = verrouillage temporaire après un virage
    \item On interdit aussi les demi-tours inutiles
\end{itemize}
\end{ideesbox}

\subsection{Détection des collisions \timing{15 sec}}

\begin{ideesbox}
\begin{itemize}[noitemsep]
    \item Comparaison des rectangles de chaque véhicule
    \item Tolérance de 2 pixels pour éviter les faux positifs
\end{itemize}
\end{ideesbox}

\subsection{Encodage UTF-8 \timing{15 sec}}

\begin{ideesbox}
\begin{itemize}[noitemsep]
    \item Le plan utilise des caractères Unicode (flèches, bordures)
    \item Nécessite \code{ncursesw} et \code{setlocale()}
\end{itemize}
\end{ideesbox}

\textbf{Transition :} \textit{"Pour conclure..."}

%==============================================================================
\section{Conclusion}
%==============================================================================

\begin{timingbox}[30 secondes]
\end{timingbox}

\begin{ideesbox}
\begin{itemize}[noitemsep]
    \item Projet formateur : gestion mémoire, structures de données, ncurses
    \item Résultat fonctionnel : navigation fluide et stable
    \item Améliorations possibles : types de véhicules, sauvegarde, multijoueur
\end{itemize}
\end{ideesbox}

\begin{phrasebox}
"Ce projet nous a permis de mettre en pratique de nombreux concepts de programmation C. Le résultat fonctionne bien et nous sommes satisfaits du comportement des véhicules. Merci pour votre attention, nous sommes prêts à répondre à vos questions."
\end{phrasebox}

\newpage

%##############################################################################
\part{Questions de Jury}
\addcontentsline{toc}{part}{Partie II : Questions de Jury (5 minutes)}
%##############################################################################

\begin{tcolorbox}[
    enhanced,
    colback=accent!5,
    colframe=accent,
    arc=5pt,
    boxrule=2pt,
    title={\color{white}\faIcon{exclamation-circle} Note importante},
    fonttitle=\bfseries,
    attach boxed title to top center={yshift=-3mm},
    boxed title style={colback=accent, arc=3pt}
]
Ces questions sont basées sur une \textbf{lecture réelle du code} et du comportement du programme. Elles correspondent à ce qu'un jury technique pourrait demander en observant le projet.

\textbf{Conseil :} Ne pas réciter les réponses par cœur, mais comprendre les concepts pour pouvoir reformuler avec ses propres mots.
\end{tcolorbox}

\vspace{0.5cm}

%==============================================================================
\section{Fonctionnement du code}
%==============================================================================

\begin{questionbox}[Q1 : Comment se déroule un tick de simulation ?]
\end{questionbox}

\begin{reponsebox}
À chaque tick (environ 57ms), on fait dans l'ordre :
\begin{enumerate}[noitemsep, leftmargin=1.5cm]
    \item[\faIcon{keyboard}] Lecture des touches clavier (E, S, Q)
    \item[\faIcon{plus-circle}] Gestion du spawn : ajout de véhicules à la file selon un cooldown
    \item[\faIcon{door-open}] Transfert d'un véhicule de la file vers le parking (si barrière ouverte)
    \item[\faIcon{arrows-alt}] Déplacement de tous les véhicules actifs
    \item[\faIcon{car-crash}] Détection des collisions
    \item[\faIcon{hourglass-end}] Vérification des timeouts en file d'attente
    \item[\faIcon{desktop}] Rafraîchissement de l'affichage
\end{enumerate}
Le tout est dans la fonction \code{executer\_boucle\_jeu()} de \code{jeu.c}.
\end{reponsebox}

\vspace{0.3cm}

\begin{questionbox}[Q2 : Comment une voiture décide-t-elle de sa direction ?]
\end{questionbox}

\begin{reponsebox}
On utilise un système de \keyword{scoring} dans \code{choisir\_direction\_stable()} :
\begin{itemize}[noitemsep]
    \item On compte les cases libres dans chaque direction (N, S, E, O)
    \item On compte les flèches de circulation dans la zone
    \item On ajoute des bonus vers la cible (place ou sortie)
    \item On choisit la direction avec le meilleur score
\end{itemize}
Mais si le score de la direction actuelle est proche du meilleur (à 2 points près), on garde la direction actuelle pour éviter les oscillations. C'est l'\keyword{hystérésis}.
\end{reponsebox}

\vspace{0.3cm}

\begin{questionbox}[Q3 : Que se passe-t-il si une voiture ne peut pas avancer ?]
\end{questionbox}

\begin{reponsebox}
Dans \code{deplacer\_vehicule()}, si la direction actuelle est bloquée :
\begin{enumerate}[noitemsep]
    \item On teste les 4 directions dans l'ordre N, S, E, O
    \item Si une direction est libre, on la prend et on met à jour le sprite
    \item Si aucune direction n'est possible, le véhicule attend le prochain tick
\end{enumerate}
C'est une solution simple mais ça peut créer des embouteillages si le parking est trop chargé.
\end{reponsebox}

\vspace{0.3cm}

\begin{questionbox}[Q4 : Comment est gérée la suppression d'un véhicule ?]
\end{questionbox}

\begin{reponsebox}
Quand un véhicule atteint la sortie (vérifié par \code{vehicule\_a\_sortie()}) :
\begin{enumerate}[noitemsep]
    \item On l'ajoute à un tableau temporaire \code{vehicules\_a\_supprimer}
    \item Après avoir traité tous les véhicules, on appelle \code{detruire\_vehicule\_specifique()}
    \item Cette fonction parcourt la liste chaînée, met à jour les pointeurs, et libère la mémoire
\end{enumerate}
On ne supprime pas pendant le parcours pour éviter de corrompre la liste chaînée.
\end{reponsebox}

\vspace{0.3cm}

\begin{questionbox}[Q5 : À quoi sert le global\_frame\_counter ?]
\end{questionbox}

\begin{reponsebox}
C'est un compteur global qui s'incrémente à chaque tick. On l'utilise pour :
\begin{itemize}[noitemsep]
    \item Calculer le temps de stationnement (frame actuelle - frame d'arrivée)
    \item Calculer le timeout en file d'attente
    \item Déterminer quand un véhicule garé doit repartir (300-600 frames)
\end{itemize}
C'est plus fiable que \code{time()} car ça dépend du framerate du jeu, pas du temps réel.
\end{reponsebox}

%==============================================================================
\section{Stabilité}
%==============================================================================

\begin{questionbox}[Q6 : Comment évitez-vous les oscillations de direction ?]
\end{questionbox}

\begin{reponsebox}
On a mis en place plusieurs mécanismes :
\begin{description}[leftmargin=0.5cm, style=nextline]
    \item[\textcolor{primary}{\faIcon{lock} Lane-lock}] Après un virage, le véhicule est "verrouillé" pendant 5 frames (constante \code{TURN\_LOCK\_DURATION})
    \item[\textcolor{secondary}{\faIcon{balance-scale} Hystérésis}] Si le score de la direction actuelle est à moins de 2 points du meilleur, on garde la direction actuelle
    \item[\textcolor{success}{\faIcon{road} Lane-keeping}] Les changements latéraux sont bloqués sauf aux intersections ou près d'une place libre
\end{description}
\end{reponsebox}

\vspace{0.3cm}

\begin{questionbox}[Q7 : Qu'est-ce qui empêche deux voitures de se superposer ?]
\end{questionbox}

\begin{reponsebox}
Deux choses :
\begin{enumerate}[noitemsep]
    \item \textbf{Prévention} : \code{voie\_libre\_direction()} vérifie qu'il n'y a pas de véhicule dans un rectangle de 8 cases dans la direction cible avant de changer de direction
    \item \textbf{Détection} : À chaque tick, on compare tous les véhicules deux à deux avec \code{vehicules\_en\_collision()}. Si collision détectée = Game Over.
\end{enumerate}
La première empêche la plupart des collisions, la deuxième est un filet de sécurité.
\end{reponsebox}

\vspace{0.3cm}

\begin{questionbox}[Q8 : Que se passe-t-il en cas de bouchon à l'entrée ?]
\end{questionbox}

\begin{reponsebox}
Si la file d'attente est pleine ou si l'entrée est bloquée :
\begin{itemize}[noitemsep]
    \item Les nouveaux véhicules ne sont pas créés (capacité max de 10)
    \item Les véhicules en attente ont un \keyword{timeout} (30s en Normal, 20s en Hard)
    \item Si le timeout est dépassé, le véhicule disparaît et on perd des points (pénalité)
\end{itemize}
C'est le mécanisme de \code{traiter\_timeout\_attente()} dans \code{jeu.c}.
\end{reponsebox}

\vspace{0.3cm}

\begin{questionbox}[Q9 : Pourquoi utiliser une tolérance de 2 pixels pour les collisions ?]
\end{questionbox}

\begin{reponsebox}
Sans tolérance, on avait beaucoup de \keyword{faux positifs} :
\begin{itemize}[noitemsep]
    \item Deux véhicules qui se croisent de près déclenchaient une collision même sans contact réel
    \item Les sprites des véhicules ont des "bords vides" (espaces dans le sprite)
\end{itemize}
La tolérance de 2 réduit la hitbox effective, ce qui donne un comportement plus naturel et évite les Game Over injustes.
\end{reponsebox}

%==============================================================================
\section{Choix de conception}
%==============================================================================

\begin{questionbox}[Q10 : Pourquoi téléporter les véhicules vers les places au lieu de les faire se garer progressivement ?]
\end{questionbox}

\begin{reponsebox}
On a essayé le parking progressif mais c'était très compliqué :
\begin{itemize}[noitemsep]
    \item Il fallait gérer des manœuvres (marche arrière, créneaux)
    \item Les collisions avec les voitures déjà garées étaient fréquentes
    \item Le véhicule pouvait rester bloqué à mi-chemin
\end{itemize}
La \keyword{téléportation} est simple et fiable : dès qu'on détecte une flèche de parking et une place libre sur la rangée, on place directement le véhicule. C'est moins réaliste mais ça fonctionne bien.
\end{reponsebox}

\vspace{0.3cm}

\begin{questionbox}[Q11 : Pourquoi avoir séparé collision.c du reste de mouvement.c ?]
\end{questionbox}

\begin{reponsebox}
Pour la \keyword{lisibilité} et la \keyword{maintenance} :
\begin{itemize}[noitemsep]
    \item \code{mouvement.c} fait déjà plus de 1000 lignes, on ne voulait pas tout mélanger
    \item La détection de collision est une logique à part (tests de rectangles)
    \item Si on veut changer l'algorithme de collision, on touche qu'un seul fichier
\end{itemize}
C'est le principe de \keyword{séparation des responsabilités}.
\end{reponsebox}

\vspace{0.3cm}

\begin{questionbox}[Q12 : Pourquoi utiliser des wchar\_t pour le plan au lieu de char ?]
\end{questionbox}

\begin{reponsebox}
Le plan utilise des caractères \keyword{Unicode} comme les flèches ($\leftarrow$, $\rightarrow$, $\uparrow$, $\downarrow$) et les bordures. Ces caractères prennent 3 octets en UTF-8.

Avec des \code{char}, on aurait dû manipuler des séquences multi-octets, c'est pénible. Avec \code{wchar\_t}, chaque case du tableau = un caractère, c'est plus simple pour les comparaisons et les tests.
\end{reponsebox}

\vspace{0.3cm}

\begin{questionbox}[Q13 : Pourquoi stocker les places dans un tableau fixe de 50 ?]
\end{questionbox}

\begin{reponsebox}
C'est un choix de \keyword{simplicité} :
\begin{itemize}[noitemsep]
    \item On sait que notre plan n'aura jamais plus de 50 places
    \item Pas besoin de gérer l'allocation dynamique et le redimensionnement
    \item L'accès par index est direct et rapide
\end{itemize}
Si on voulait supporter des plans plus grands, on pourrait passer à un tableau dynamique, mais pour notre projet c'était suffisant.
\end{reponsebox}

%==============================================================================
\section{Lecture de code}
%==============================================================================

\begin{questionbox}[Q14 : À quoi sert le champ lane\_lock\_counter dans VehiculeState ?]
\end{questionbox}

\begin{reponsebox}
C'est le compteur du mécanisme de \keyword{lane-lock} :
\begin{itemize}[noitemsep]
    \item Quand un véhicule tourne, on met \code{lane\_lock\_counter = 5}
    \item À chaque tick, on décrémente ce compteur
    \item Tant qu'il est $> 0$, le véhicule ne peut pas rechanger de direction
\end{itemize}
Ça évite les oscillations où le véhicule changerait de direction à chaque frame.
\end{reponsebox}

\vspace{0.3cm}

\begin{questionbox}[Q15 : Pourquoi get\_state() utilise un tableau statique au lieu d'une map ?]
\end{questionbox}

\begin{reponsebox}
En C, on n'a pas de \keyword{map native} comme en C++ ou Python. On a fait simple :
\begin{itemize}[noitemsep]
    \item Un tableau statique de 20 slots (constante \code{MAX\_VOITURES})
    \item On cherche le véhicule par parcours linéaire
    \item Si pas trouvé, on prend le premier slot libre
\end{itemize}
Avec 20 véhicules max, le parcours linéaire est assez rapide. Une vraie hashmap aurait été overkill pour ce cas d'usage.
\end{reponsebox}

\vspace{0.3cm}

\begin{questionbox}[Q16 : Que se passerait-il si on enlevait setlocale() au début de main() ?]
\end{questionbox}

\begin{reponsebox}
Les caractères Unicode s'afficheraient mal :
\begin{itemize}[noitemsep]
    \item Les flèches ($\leftarrow$, $\rightarrow$) deviendraient des "?" ou des caractères bizarres
    \item Les bordures du plan seraient illisibles
    \item Le programme fonctionnerait mais l'affichage serait cassé
\end{itemize}
\code{setlocale(LC\_ALL, "")} dit au programme d'utiliser l'encodage du terminal (généralement UTF-8).
\end{reponsebox}

\vspace{0.3cm}

\begin{questionbox}[Q17 : Pourquoi deplacer\_tous\_vehicules() retourne 1 en cas de collision au lieu de gérer le game over directement ?]
\end{questionbox}

\begin{reponsebox}
C'est pour \keyword{séparer les responsabilités} :
\begin{itemize}[noitemsep]
    \item \code{mouvement.c} gère le déplacement, il ne connaît pas l'affichage
    \item \code{jeu.c} gère la logique du jeu et l'affichage du game over
    \item En retournant juste un booléen, \code{mouvement.c} reste indépendant
\end{itemize}
C'est \code{gerer\_collisions\_et\_game\_over()} dans \code{jeu.c} qui affiche l'écran de game over.
\end{reponsebox}

%==============================================================================
\section{Questions binôme}
%==============================================================================

\begin{questionbox}[Q18 : Comment avez-vous réparti le travail ?]
\end{questionbox}

\begin{reponsebox}
\textit{(À adapter selon votre répartition réelle)}

On a divisé par modules :
\begin{itemize}[noitemsep]
    \item \textbf{Personne A} : Affichage ncurses, chargement du plan, structures de données
    \item \textbf{Personne B} : Déplacement des véhicules, détection de collision, boucle de jeu
\end{itemize}
On a travaillé ensemble sur les parties communes (intégration, débogage). On faisait des points réguliers pour synchroniser nos avancées.
\end{reponsebox}

\vspace{0.3cm}

\begin{questionbox}[Q19 : Comment avez-vous testé le projet ?]
\end{questionbox}

\begin{reponsebox}
Plusieurs méthodes :
\begin{itemize}[noitemsep]
    \item[\faIcon{gamepad}] \textbf{Tests manuels} : On lançait le jeu et on observait le comportement
    \item[\faIcon{vial}] \textbf{Tests unitaires} : Tests Unity pour la liste chaînée (\code{test\_liste.c})
    \item[\faIcon{globe}] \textbf{Tests d'encodage} : Test spécifique pour les caractères UTF-8
    \item[\faIcon{memory}] \textbf{Valgrind} : Vérification des fuites mémoire
\end{itemize}
\end{reponsebox}

\vspace{0.3cm}

\begin{questionbox}[Q20 : Quel a été le bug le plus difficile à corriger ?]
\end{questionbox}

\begin{reponsebox}
\textit{(À adapter selon votre expérience réelle)}

Les \keyword{oscillations des véhicules}. Au début, une voiture changeait de direction à chaque frame et faisait du "zigzag". On a dû :
\begin{enumerate}[noitemsep]
    \item Comprendre pourquoi les scores changeaient à chaque tick
    \item Implémenter le lane-lock
    \item Ajouter l'hystérésis de direction
    \item Bloquer les changements latéraux hors intersections
\end{enumerate}
Ça nous a pris plusieurs jours pour stabiliser tout ça.
\end{reponsebox}

%------------------------------------------------------------------------------
% PAGE FINALE
%------------------------------------------------------------------------------
\newpage

\begin{center}
\vspace*{3cm}

\begin{tikzpicture}
    % Fond
    \fill[primary!10, rounded corners=20pt] (-7,-4) rectangle (7,4);

    % Texte principal
    \node[font=\Huge\bfseries, color=primary] at (0,2) {\faIcon{graduation-cap} Bonne chance !};

    % Ligne décorative
    \draw[secondary, line width=2pt] (-4,0.8) -- (4,0.8);

    % Conseils
    \node[font=\large, color=dark, align=center] at (0,-0.5) {
        Restez calmes, répondez honnêtement,\\[0.3cm]
        et n'ayez pas peur de dire\\[0.3cm]
        \textbf{"Je ne sais pas"} si vous ne connaissez pas la réponse.
    };

    % Icônes en bas
    \node at (-3,-3) {\textcolor{success}{\fontsize{30}{30}\selectfont\faIcon{smile}}};
    \node at (0,-3) {\textcolor{warning}{\fontsize{30}{30}\selectfont\faIcon{lightbulb}}};
    \node at (3,-3) {\textcolor{secondary}{\fontsize{30}{30}\selectfont\faIcon{thumbs-up}}};
\end{tikzpicture}

\vspace{2cm}

\textcolor{dark}{\large\faIcon{quote-left} \textit{Une réponse honnête vaut mieux qu'une réponse inventée} \faIcon{quote-right}}

\end{center}

\end{document}
